\documentclass[a4paper,14pt]{extarticle}
\usepackage{array,amsmath,fancybox,amsfonts,graphicx,amsthm,amssymb,wasysym,latexsym,slashed,anysize,subfigure,calc}
\usepackage[spanish]{babel}
\usepackage[utf8]{inputenc}
\usepackage{fancyhdr,float,hyperref}
\usepackage{lettrine}


%lewis y química
\usepackage{chemformula,elements}
%chemformula, provides \ch{} with the !(<below>)(<formula>) syntax and \chlewis;
%elements, provides \elconf and \writeelconf.
%ex:
%\ch{
%  !(\elconf{Li})( "\chlewis{180.}{Li}" ) +
%  !(\elconf{F})( "\chlewis{0.90:180:270:}{F}" )
%   ->
%  !(\writeelconf{2})( Li+ ) +
%  !(\writeelconf{2,2+6})( "\chlewis{0:90:180:270:}{F}" {}- )
%}


%entorno para química: texto y math mode
%\newcommand*\chem[1]{\textup{#1}}
\newcommand*\chem[1]{\ensuremath{\mathrm{#1}}}

%column
\usepackage{multicol}
%uso: \begin{multicols}{3}

%enumitem: configurar listas
\usepackage{enumitem}

%para los entornos de soluciones
\usepackage{mdframed,xcolor,color}
	%uso: 	entorno (begin,end): {mdframed}[backgroundcolor=brown!10] 

%subrayado y highlight. Tiene que estar detrás de usepackage[color]
\usepackage{soul}
	%\textst{normal} 
	\DeclareRobustCommand{\hlcyan}[1]{{\sethlcolor{cyan}\hl{#1}}} %en cyan!
		%\hl{highlighted text} %uso


%tamaños fuentes extra: 8,9,14,17,20
	\usepackage{extsizes}
	
 	%fuente librebaskerville
 	\usepackage[T1]{fontenc}
	\usepackage{librebaskerville}
	
	%fracciones $$ en texto más elaboradas
	 \usepackage{xfrac}
	 
	%permitir que una ecuación larga esté separada por un salto de página
	\allowdisplaybreaks[4]
		
	%gráficos varios 
	\usepackage{graphics}
	
	%fancy tables
	\usepackage{booktabs}
	\usepackage[small,centerlast,up]{caption}

	
	%mejora en la exposición de guarismos
	\usepackage{numprint}
	
\usepackage{everypage,draftwatermark} %Marca de agua
\SetWatermarkText{\vspace*{2cm}

\hspace*{2cm}\parbox{2\textwidth}{\centering DIST.\\B1-FyQ}}
\SetWatermarkScale{1.5}
\SetWatermarkLightness{0.9}

\DeclareMathOperator{\lin}{lin}
\DeclareMathOperator{\id}{\textbf{id}}
\DeclareMathOperator{\modulo}{mod}

\newcommand{\mc}[3]{\multicolumn{#1}{#2}{#3}}
\newcommand{\halmos}{
\vspace*{-3ex}

\begin{flushright}
\rule{1mm}{2.5mm} \end{flushright}
}
\renewcommand{\qed}{\halmos}


%versores: %uso: \uveci\ne\hat{\imath}_{\uvec{i}}+\uvec{j}+\uvec{k}$
\usepackage{bm}
\newcommand{\uveci}{{\bm{\hat{\textnormal{\bfseries\i}}}}}
\newcommand{\uvecj}{{\bm{\hat{\textnormal{\bfseries\j}}}}}
\DeclareRobustCommand{\uvec}[1]{{%
  \ifcsname uvec#1\endcsname
     \csname uvec#1\endcsname
   \else
    \bm{\hat{\mathbf{#1}}}%
   \fi
}}


%algún nuevo operador
\DeclareMathOperator{\cotan}{cotan}


%para textbf en entorno enumerate
\renewcommand{\labelenumi}{%
 \textbf{\theenumi}.
}

%a4
\marginsize{1cm}{1cm}{1cm}{1cm}%requiere el paquete anysize,izq,dcha,arriba,abajo

%a5
%\marginsize{1cm}{1cm}{1cm}{.5cm}%requiere el paquete anysize,izq,dcha,arriba,abajo

\newtheorem{teo}{{\sc \textbf{Teorema}}}
\newtheorem{defi}{{\sc \textbf{Definición}}}
\newtheorem{coro}{{\sc \textbf{Corolario}}}


%opening
\title{\textbf{	}}
\date{}
\pagestyle{fancy}
\fancyhead[LO,RE]{{\small DOSSIER DE PROBLEMAS}}
\fancyhead[RO,LE]{{\small B1-DIST. FÍSICA Y QUÍMICA}}
%\fancyhead[CO,CE]{Ricardo Torres\\http://tinyurl.com/iesCMGfyq}



\begin{document}
\pagenumbering{gobble}

%	\section*{Ejercicios propuestos Tema 3: óptica ondulatoria}

%\begin{center}
%	\hl{2, 4, 10 p\'ag. 241}\hspace{2cm}\hl{24 p\'ag. 242}
%	\end{center}

%\vspace*{.1cm}
%
%	\section*{Ejercicios adicionales sobre la luz}
%


\begin{enumerate}
  \setlength\itemsep{1.5em}
    
    
    \item \textbf{(2 PUNTOS)} 
    
    Formula y nombra:
    	%\begin{multicols}{2}
    		\begin{enumerate}[label=\alph*)]
    			\item Los ácidos ternarios de los tres primeros elementos del grupo 15.
    			\item Las sales que forma el azufre con el segundo y tercer elemento del grupo cuyo último electrón corresponda a un estado $ns^2$.
    			\item Los peróxidos que se forman con los metales alcalinos correspondientes a los periodos $3$ y $4$. 
    			\item Las sales ternarias que forma el hierro con los aniones resultantes de la disociación completa de los ácidos ternarios correspondientes a los elementos con configuración $\elconf{se}$ y $\elconf{i}$.
    		\end{enumerate}
    		\qed
    	%\end{multicols}
    	
    \item \textbf{(1 PUNTO)}
    
    	\begin{enumerate}[label=\alph*)]
    		\item ¿Cuál es la diferencia entre un alcano y un alquino?
    		\item  Nombra y formula (en forma semidesarrollada) los cinco primeros alquenos en \textit{todas} las formas \textit{distintas} en que sea posible que contengan dobles enlaces.
    		\item ¿Por qué el grupo característico de las cetonas no puede ir asociado a los carbonos en los extremos de una cadena lineal?
    	\end{enumerate} 
    \qed
    
    
    \item \textbf{(1 PUNTO)}
    
     A la vista del proceso que hemos analizado para la formación de sales ternarias a partir la disociación de ácidos oxoácidos, ¿qué relación existe entre un ácido carboxílico y un éster? 
    \qed
    
    
    \item \textbf{(3 PUNTOS)} 
    
    Considera los elementos desconocidos $^{8}\alpha$, $^{18}\beta$, $^{34}\delta$ y $^{38}\epsilon.$ 
    	\begin{enumerate}[label=\alph*)]
    		\item Escribe las configuraciones electrónicas de los cuatro elementos. 
    		\item Justifica el tipo de enlace y la fórmula del compuesto que forman los elementos $\delta$ y $\epsilon$.
    		\item Justifica la existencia de la molécula $\alpha_2$ e indica el tipo de enlace que se produce.
    		\item Indica el tipo de enlace que formarán los elementos $\beta$ y $\alpha$.
    	\end{enumerate}
    \qed
    	
    
    \newpage
    
    \item \textbf{(3 PUNTOS)} 
    
    El modelo de Bohr para el hidrógeno permite explicar la existencia de un espectro atómico discreto debido a un mecanismo de transición entre órbitas. En efecto, el único electrón del hidrógeno es capaz de pasar de una órbita más externa a una más interna emitiendo luz; o pasar de una órbita más interna a una más externa absorbiendo luz.
    
    La energía de un electrón en una órbita descrita por el número cuántico $n$ puede escribirse como
    	$$E_n=-\frac{13.6}{n^2},$$
    medida en $eV$, una unidad de energía que equivale aproximadamente a $1.6\cdot 10^{-19}\;J$.
    		
    	\begin{enumerate}[label=\alph*)]
    		\item La cantidad de energía luminosa emitida o absorbida en la transición de un electrón desde un nivel $n_1$ a un nivel $n_2$ es $$\Delta E=|E_{n_1}-E_{n_2}|.$$
    		
    	 Calcula las cuatro energías, en $eV$, que puede emitir un electrón cuando \textit{cae} al nivel fundamental $n=1$ procedente de las órbitas correspondientes a $n=2$, $n=3$, $n=4$ y $n=5$.
    		\item Planck propuso que la energía luminosa emitida o absorbida en una transición correspondía a la emisión o absorción de una partícula elemental llamada fotón que podía caracterizarse por una frecuencia $f$, medida en $s^{-1}$, según la expresión 
    		$$\Delta E=hf,$$
    		donde $h\simeq 6.63\cdot 10^{-34}\;J\cdot s$ es la llamada constante de Planck. 
    		Calcula la frecuencia asociada a los fotones correspondientes a las transiciones del apartado a).
    		\item Años antes Rydberg, de forma completamente experimental, había llegado a la conclusión de que la frecuencia de un fotón emitido como resultado de la transición de un electrón desde el nivel $n_i$ a un nivel inferior $n_f$ podía escribirse como 
    		$$f=\beta\left(\frac{1}{n_f^2}-\frac{1}{n_i^2}\right),$$
    		donde $\beta$ es una cierta constante. Calcula el valor de $\beta$ y sus unidades a partir de la información de los apartados a) y b).
    	\end{enumerate}
		\qed    	

\end{enumerate}

\end{document}

